\chapter{Modèles réactifs à la file et actifs à gros tick}\label{ch:hf:queue-reactive-models-and-large-tick-assets}

Sur un titre dont la fourchette vaut un tick presque toute la journée, chaque teneur de marché affiche les mêmes deux prix. Ce qui les sépare, c'est la place qu'ils occupent dans la
file, et le moment où ils la quittent. Sur le titre simulé de ce chapitre, la fourchette vaut un tick 96~\% du temps. Un teneur de marché à un lot qui entre toujours dans les
meilleures files gagne 0.03 tick par action sur les dix secondes suivant ses exécutions~; un autre qui n'entre que du côté dont la file représente au moins 30~\% des deux,
et quitte la file lorsque cette part baisse, gagne 0.39 sur un tiers du volume. La même règle d'entrée sans la règle de sortie perd 0.12. Quitter la file est la partie qui
rapporte.

\section{Le régime à gros tick}

Le livre~10, chapitre~6 classe les instruments selon leur pas de cotation relatif~: dans un \emph{actif à gros tick} la fourchette vaut presque toujours un tick
et les meilleures files sont longues, alors que dans un \emph{actif à petit tick} la fourchette couvre de nombreux ticks et les files sont courtes. Cette distinction détermine ce que le
teneur de marché contrôle. Dans un actif à petit tick, la distance de la cotation au mid est un choix continu, la profondeur du chapitre~3. Dans un actif à gros tick, le
prix est fixé par la grille des ticks, et les seuls choix sont d'être dans une meilleure file ou non, et où s'y placer.

Le titre simulé du livre~7 est à gros tick. Sur une heure, sa fourchette vaut un tick 96.3~\% du temps et la médiane de la meilleure file compte 17 lots de 100
actions. Le teneur de marché explorateur du chapitre~3 n'a obtenu aucune exécution à un tick derrière le meilleur prix. Tout ce qu'un teneur de marché y gagne, il le gagne dans les deux
meilleures files.

\section{Intensités réactives à la file}

\begin{definition}[Modèle réactif à la file]\label{def:hf:queue-reactive-models-and-large-tick-assets:qr}
Un \emph{modèle réactif à la file}\index{modèle réactif à la file} (queue-reactive model) décrit chaque meilleure file d'un carnet d'ordres comme un processus de naissance et de mort dont les taux dépendent de la
taille courante $q$ de la file~: les ordres à cours limité rejoignent son arrière au taux $\lambda^L(q)$, les annulations retirent des ordres au repos au taux $\lambda^C(q)$, les ordres au marché
retirent des ordres à son front au taux $\lambda^M(q)$. Quand une file se vide, le prix se déplace d'un tick et de nouvelles files sont tirées à partir
d'une loi de régénération.
\end{definition}

Huang, Lehalle et Rosenbaum ont introduit le modèle à partir de données réelles de carnet d'ordres~: entre deux changements d'un prix de référence, le carnet est un système de files d'attente
markovien dont les intensités ne dépendent que de son état courant. C'est la chaîne de naissance et de mort du livre~4, chapitre~8, avec des taux lus dans les
données au lieu d'être supposés constants, comme dans le modèle antérieur de Cont, Stoikov et Talreja.

\begin{method}[Estimation des intensités]\label{met:hf:queue-reactive-models-and-large-tick-assets:estimate}
Rejouer les messages ordre par ordre. Pour chaque message, noter la taille de chaque meilleure file juste avant celui-ci et le temps écoulé depuis le message précédent~;
imputer ce temps à la taille. Compter les lots ajoutés, annulés et exécutés contre chaque meilleure file à chaque taille. L'intensité pour la taille $q$
est ce nombre divisé par le temps passé à $q$~; regrouper les deux côtés si le carnet est symétrique.
\end{method}

La \Cref{fig:hf:queue-reactive-models-and-large-tick-assets:intensities} montre l'estimation sur une heure simulée. Les annulations croissent avec la file,
de 0.34 lot par seconde pour un lot à environ 1.5 pour dix, parce que chaque ordre au repos s'annule à son propre rythme~; les ordres au marché atteignent les petites files plus
souvent que les grandes (0.55 lot par seconde pour deux lots, environ 0.3 pour quinze). Les ordres à cours limité arrivent à un rythme de 1.3 à 2.3 lots par seconde, sans tendance liée à la taille.

\begin{omfigure}
\begin{tikzpicture}
\begin{axis}[width=11cm,height=5.2cm,xlabel={\small taille de la file avant l'événement (lots)},ylabel={\small lots par seconde},
  tick label style={font=\footnotesize},xmin=1,xmax=25,ymin=0,ymax=2.6,grid=major,grid style={gray!25},table/col sep=comma,
  legend style={font=\scriptsize,at={(0.5,-0.3)},anchor=north},legend columns=3]
\addplot[omDef,thick,mark=*,mark size=1pt] table[x=q,y=L]{figdata/hft/05-queue-reactive-models-and-large-tick-assets/intensities.csv};
\addplot[omThm,thick,dashed,mark=square*,mark size=1pt] table[x=q,y=C]{figdata/hft/05-queue-reactive-models-and-large-tick-assets/intensities.csv};
\addplot[black,thick,mark=triangle*,mark size=1.2pt] table[x=q,y=M]{figdata/hft/05-queue-reactive-models-and-large-tick-assets/intensities.csv};
\legend{ordres à cours limité $\lambda^L$,annulations $\lambda^C$,ordres au marché $\lambda^M$}
\end{axis}
\end{tikzpicture}

\omcaption{Intensités réactives à la file des meilleures files du titre simulé, regroupées sur les côtés cours acheteur et cours vendeur, selon la taille de la file juste avant chaque
événement~; une heure simulée. Données~: \texttt{firm.qreactive.estimate} sur \texttt{firm.tape}.}\label{fig:hf:queue-reactive-models-and-large-tick-assets:intensities}
\end{omfigure}

\section{La valeur d'une place dans la file}

Le livre~10, chapitre~6 définit la \emph{valeur de la position dans la file} d'un ordre au repos~: ce que l'ordre devrait rapporter depuis sa place actuelle jusqu'à ce
qu'il soit exécuté ou annulé. Moallemi et Yuan la décomposent en une part statique, la fourchette gagnée face aux coûts de sélection adverse qui croissent avec la
position dans la file, et une part dynamique, l'option de quitter la file plus tard~; pour certains actifs à gros tick, ils l'ont trouvée du même ordre de grandeur que la fourchette.

Le \omterm{def:hf:queue-reactive-models-and-large-tick-assets:qr}{modèle réactif à la file} calibré le calcule par simulation. On place un lot des nôtres dans la file acheteuse avec $n$ lots devant lui~; la file acheteuse compte $Q^b$
lots au total et la file vendeuse $Q^a$. Les ventes au marché consomment la file acheteuse par l'avant, les annulations touchent au hasard les autres ordres, les ordres à cours limité rejoignent la
file par l'arrière. L'ordre est exécuté lorsqu'une vente au marché l'atteint~; il est abandonné, de valeur nulle, si la file vendeuse se vide en premier (le prix monte, s'éloignant
de lui) ou après soixante secondes. Après l'exécution, le mid est suivi pendant dix secondes, chaque file vidée le déplaçant d'un tick, et l'ordre est évalué
à ce moment-là~: un demi-tick plus le déplacement du mid (le \cref{lst:hf:queue-reactive-models-and-large-tick-assets:mc}).

\begin{omfigure}
\begin{tikzpicture}
\begin{axis}[width=11cm,height=5.2cm,xlabel={\small taille de la file de notre côté $Q^b$ (lots)},ylabel={\small valeur (ticks par action)},xmode=log,
  log basis x=10,xtick={2,5,10,20},xticklabels={2,5,10,20},xmin=1.8,xmax=22,ymin=0,ymax=0.7,tick label style={font=\footnotesize},
  grid=major,grid style={gray!25},table/col sep=comma,legend style={font=\scriptsize,at={(0.5,-0.3)},anchor=north},legend columns=2]
\addplot[omDef,thick,mark=*,mark size=1.2pt] table[x=same,y=front2]{figdata/hft/05-queue-reactive-models-and-large-tick-assets/values.csv};
\addplot[omDef,thick,dashed,mark=o,mark size=1.2pt] table[x=same,y=back2]{figdata/hft/05-queue-reactive-models-and-large-tick-assets/values.csv};
\addplot[omThm,thick,mark=square*,mark size=1.2pt] table[x=same,y=front20]{figdata/hft/05-queue-reactive-models-and-large-tick-assets/values.csv};
\addplot[omThm,thick,dashed,mark=square,mark size=1.2pt] table[x=same,y=back20]{figdata/hft/05-queue-reactive-models-and-large-tick-assets/values.csv};
\legend{{avant, $Q^a=2$},{arrière, $Q^a=2$},{avant, $Q^a=20$},{arrière, $Q^a=20$}}
\end{axis}
\end{tikzpicture}

\omcaption{Valeur d'un lot au repos dans la file acheteuse, à son avant et à son arrière, dans le \omterm{def:hf:queue-reactive-models-and-large-tick-assets:qr}{modèle réactif à la file} calibré, selon la taille de la file acheteuse,
pour une file vendeuse courte et une file vendeuse longue~: P\&L espéré par action, évalué dix secondes après l'exécution (nul si l'ordre n'est pas exécuté dans les soixante secondes ou si la
file vendeuse se vide en premier)~; 4\,000 trajectoires simulées par point. Données~: \texttt{hf\_queues.values}.}\label{fig:hf:queue-reactive-models-and-large-tick-assets:values}
\end{omfigure}

Trois régularités apparaissent dans la \cref{fig:hf:queue-reactive-models-and-large-tick-assets:values}. L'avant vaut davantage que l'arrière lorsque notre
file est longue~: à dix lots avec une file vendeuse courte, 0.60 contre 0.32 tick. Lorsque notre file est courte et l'autre longue, l'avant et l'arrière valent
peu, et presque identiques (0.12 pour deux lots contre vingt)~: une exécution est presque certaine, car la file est sur le point d'être consommée, et le prix suit.
Lorsque l'autre file est courte, un ordre à l'arrière risque de ne jamais être exécuté, car le prix s'éloigne d'abord.

\begin{remark}[Ce que le modèle ne voit pas]\label{rem:hf:queue-reactive-models-and-large-tick-assets:info}
Chaque valeur de la figure est positive, pourtant le teneur de marché du chapitre~1, toujours dans ces files, a perdu sa fourchette face à la sélection adverse. Le modèle calibré
ne connaît que les tailles des files. Sur le marché simulé, comme sur les marchés réels, certains ordres au marché proviennent de traders qui savent où se situe le prix efficient,
et ils arrivent lorsque les cotations sont périmées. Un \omterm{def:hf:queue-reactive-models-and-large-tick-assets:qr}{modèle réactif à la file} capture la course mécanique pour l'avant~; l'information contenue dans le flux doit être
mesurée sur ses propres exécutions, par des mark-outs, et ajoutée à celui-ci (chapitre~6).
\end{remark}

\section{Quand entrer dans la file et quand la quitter}

\begin{definition}[Règle d'entrée dans la file, règle de sortie de la file]\label{def:hf:queue-reactive-models-and-large-tick-assets:rules}
Une \emph{règle d'entrée dans la file}\index{règle d'entrée dans la file} (queue-join rule) décide s'il faut placer un ordre à l'arrière d'une meilleure file, étant donné l'état du carnet d'ordres~; une
\emph{règle de sortie de la file}\index{règle de sortie de la file} (queue-exit rule) décide s'il faut annuler un ordre déjà au repos, étant donné l'état du carnet et la position de l'ordre
dans sa file.
\end{definition}

Les deux règles sont des décisions différentes~: entrer dans une file n'abandonne rien, la quitter fait perdre sa place. Les règles de ce chapitre utilisent un seul nombre, la part de la meilleure
file de notre côté dans les deux meilleures files, en ne comptant que les ordres des autres, $s^b=Q^b/(Q^b+Q^a)$ pour le cours acheteur (et $1-s^b$ pour le cours vendeur). Selon le modèle, un côté
dont la part est faible va être consommé~: les exécutions y sont presque certaines et suivies d'un mouvement défavorable. Un teneur de marché à un lot n'entre dans un côté
seulement lorsque sa part vaut au moins $\theta$~; la règle de sortie soit reste en place, soit quitte la file lorsque la part passe sous $\theta$, soit quitte la file sauf si moins de
deux lots sont devant (le \cref{lst:hf:queue-reactive-models-and-large-tick-assets:quoter}). Sur six séances simulées de vingt minutes~:

\begin{center}
\small
\setlength{\tabcolsep}{4pt}%
\begin{tabular}{@{}>{\raggedright\arraybackslash}p{4.2cm}>{\centering\arraybackslash}p{2.4cm}>{\centering\arraybackslash}p{2.4cm}>{\centering\arraybackslash}p{2.4cm}>{\centering\arraybackslash}p{2.4cm}@{}}
\toprule
règle ($\theta=0.3$) & actions & mark-out à dix secondes & edge à dix secondes & P\&L (écart-type)\\
& par séance & (ticks par action) & (\$ par séance) & (\$ par séance)\\
\midrule
toujours entrer & 11\,167 & 0.026 & 2.92 & 5.33 (63.67)\\
entrer, puis rester & 6\,900 & $-0.117$ & $-8.08$ & $-32.00$ (68.41)\\
entrer et quitter & 3\,400 & 0.392 & 13.33 & $-6.50$ (23.82)\\
entrer, quitter sauf en tête de file & 4\,450 & 0.032 & 1.42 & $-26.50$ (63.22)\\
\bottomrule
\end{tabular}
\end{center}

Le mark-out à dix secondes, soit la capture du spread diminuée du mouvement défavorable sur dix secondes, est l'edge de chaque exécution~; multiplié par les actions, il donne l'edge de la séance. Selon
cette mesure, c'est la règle de sortie qui fait la différence. Entrer sur une part favorable puis rester est pire que toujours entrer~: les ordres qui restent sont
précisément ceux dont le côté s'est retourné contre eux, et ce sont eux qui sont exécutés. Quitter lorsque la part baisse quadruple l'edge sur un tiers du
volume. Garder les ordres près de la tête de file, le raffinement naturel, restitue l'essentiel du gain, comme le prédisait le modèle~: lorsque l'autre côté
est sur le point de prendre la file, être en tête signifie être exécuté en premier. La colonne P\&L, qui valorise aussi l'inventaire à la fin de chaque séance, est trop
bruitée sur six séances pour classer les règles.

La \Cref{fig:hf:queue-reactive-models-and-large-tick-assets:sweep} fait varier $\theta$ pour la règle entrer et quitter. Le mark-out par action croît avec le
seuil jusqu'à 0.4~; l'edge par séance culmine à 0.3, valeur où le teneur de marché échange encore suffisamment.

\begin{omfigure}
\begin{tikzpicture}
\begin{axis}[width=11cm,height=5.2cm,axis y line*=right,axis x line=none,ylabel={\small mark-out à dix secondes (ticks par action)},
  tick label style={font=\footnotesize},xmin=-0.02,xmax=0.52,ymin=0,ymax=0.6,table/col sep=comma]
\addplot[omDef,thick,dashed,mark=square*,mark size=1.3pt] table[x=theta,y=markout10]{figdata/hft/05-queue-reactive-models-and-large-tick-assets/sweep.csv};
\label{plot:hf:qsweep}
\end{axis}
\begin{axis}[width=11cm,height=5.2cm,axis y line*=left,xlabel={\small seuil $\theta$ sur la part de notre côté dans les meilleures files},
  ylabel={\small edge à dix secondes (\$ par séance)},tick label style={font=\footnotesize},xmin=-0.02,xmax=0.52,ymin=0,ymax=15,grid=major,
  grid style={gray!25},table/col sep=comma,legend style={font=\scriptsize,at={(0.5,-0.3)},anchor=north},legend columns=2]
\addplot[omThm,thick,mark=*,mark size=1.3pt] table[x=theta,y=edge10]{figdata/hft/05-queue-reactive-models-and-large-tick-assets/sweep.csv};
\addlegendentry{edge par séance (gauche)}
\addlegendimage{/pgfplots/refstyle=plot:hf:qsweep}\addlegendentry{mark-out par action (droite)}
\end{axis}
\end{tikzpicture}

\omcaption{La règle entrer et quitter sur six séances simulées de vingt minutes~: edge à dix secondes par séance et mark-out à dix secondes par action exécutée,
en fonction du seuil $\theta$ ($\theta=0$ correspond à toujours entrer). Données~: \texttt{hf\_queues.sweep}.}\label{fig:hf:queue-reactive-models-and-large-tick-assets:sweep}
\end{omfigure}

\section{Tutoriel~: valoriser une position dans la file}\label{tut:hf:queue-reactive-models-and-large-tick-assets:tut}

\begin{tutorial}
\textbf{Objectif.} Estimer les intensités du \omterm{def:hf:queue-reactive-models-and-large-tick-assets:qr}{modèle réactif à la file} d'un carnet à gros tick, valoriser une position dans sa file et tester, sur le
marché, les règles d'entrée et de sortie. \textbf{État final~:} les \cref{fig:hf:queue-reactive-models-and-large-tick-assets:intensities,fig:hf:queue-reactive-models-and-large-tick-assets:values,fig:hf:queue-reactive-models-and-large-tick-assets:sweep}
et le tableau.
\begin{tutsteps}
\item \textbf{Estimer} avec \texttt{firm.qreactive.estimate(msgs, top, lot, qmax)} sur une heure simulée (\cref{met:hf:queue-reactive-models-and-large-tick-assets:estimate}).
\item \textbf{Valoriser la file} avec \texttt{order\_value}~: tous les états et tous les chemins avancent ensemble, un événement chacun à chaque pas. Les événements du côté acheteur
en sont le cœur.
\omcode{code/firm/qreactive/firm_qreactive.py}{121}{138}{Notre ordre dans la file acheteuse~: des entrées derrière lui, des annulations devant lui, des ordres au marché depuis la tête.}\label{lst:hf:queue-reactive-models-and-large-tick-assets:mc}
\item \textbf{Règles pour le teneur de marché.} Le harnais donne au teneur de marché la position de son ordre dans la file (\texttt{ctx.ahead}), comme le ferait un flux
ordre par ordre~; le teneur de marché lit les files des autres depuis \texttt{ctx.external}.
\omcode{code/hft/05-queue-reactive-models-and-large-tick-assets/python/hf_queues.py}{79}{92}{Entrer dans un côté selon sa part des meilleures files~; rester, quitter, ou quitter sauf en tête de file.}\label{lst:hf:queue-reactive-models-and-large-tick-assets:quoter}
\item \textbf{Comparer} avec \texttt{compare()} et \texttt{sweep()} sur des séances jamais utilisées pour l'estimation.
\end{tutsteps}
\textbf{Pour aller plus loin.} Ajouter le déséquilibre du niveau suivant~; faire dépendre le seuil de la position dans la file, en s'appuyant sur la valeur du modèle.
\end{tutorial}

\section{Construction~: le moteur réactif à la file}\label{bld:hf:queue-reactive-models-and-large-tick-assets:qreactive}

\begin{build}
\bfield{Rôle} Intensités du \omterm{def:hf:queue-reactive-models-and-large-tick-assets:qr}{modèle réactif à la file} estimées à partir de tout flux ordre par ordre, ainsi que la valeur de la position d'un ordre dans sa file, pour les décisions d'entrée et de
sortie.
\bfield{Interface} \texttt{estimate(msgs, top, lot, qmax, n\_open) -> \{q, L, C, M, time, regen\}}~; \texttt{QRModel(L, C, M, regen)}~;
\texttt{order\_value(model, ahead, same, opp, H, tmax, paths, seed) -> \{fill, value, adverse\}}~; \texttt{value\_table}. Dans
\texttt{firm.mmharness}, \texttt{ctx.ahead(cid)} pour les teneurs de marché qui définissent \texttt{wants\_queue\_position}.
\bfield{Règles} Les événements sont comptés en lots, un lot par événement. L'estimation et l'évaluation n'utilisent jamais les mêmes séances. Les valeurs sont par action, en
ticks, valorisées au mid à un horizon fixe après l'exécution.
\bfield{Tests d'acceptation} \texttt{code/firm/qreactive/tests/}~: intensités et temps passé dans chaque état calculés à la main sur un flux de quatre messages~; aucune exécution sans ordres au
marché~; la tête de file est exécutée plus souvent que l'arrière~; forme des tableaux. \texttt{code/firm/mmharness/tests/}~: la position d'un ordre dans sa file ne recule jamais.
\bfield{Pour aller plus loin} Intensités dépendant des deux files (le modèle complet)~; un terme d'information ajusté aux mark-out~; lancer le modèle sur le
flux de \texttt{firm.exchsim}.
\end{build}

\begin{omsources}
\item W.~Huang, C.-A.~Lehalle et M.~Rosenbaum, ``Simulating and analyzing order book data: the queue-reactive model'', \emph{Journal of the American
Statistical Association} 110(509), 2015.
\item C.~C.~Moallemi et K.~Yuan, ``A model for queue position valuation in a limit order book'', document de travail, 2016.
\item R.~Cont, S.~Stoikov et R.~Talreja, ``A stochastic model for order book dynamics'', \emph{Operations Research} 58(3), 2010.
\item K.~Dayri et M.~Rosenbaum, ``Large tick assets: implicit spread and optimal tick size'', \emph{Market Microstructure and Liquidity} 1(1), 2015.
\end{omsources}

\section{Exercices}

\begin{exercise}[$\star$]\label{exo:hf:queue-reactive-models-and-large-tick-assets:1}
La meilleure file est restée 200 secondes à 4 lots~; pendant ce temps, 290 lots ont été annulés de cette file. Estimer $\lambda^C(4)$.
\end{exercise}

\begin{exercise}[$\star$]\label{exo:hf:queue-reactive-models-and-large-tick-assets:2}
La file acheteuse compte 3 lots d'ordres des autres et la file vendeuse 12. Quelle est la part du côté acheteur, et la règle du chapitre, avec $\theta=0.3$, conduit-elle à entrer dans le
côté acheteur~? Et le côté vendeur~?
\end{exercise}

\begin{exercise}[$\star$]\label{exo:hf:queue-reactive-models-and-large-tick-assets:3}
Les exécutions d'un teneur de marché ont rapporté 0.392 ticks par action au bout de dix secondes, sur 3\,400 actions par séance. Avec un tick d'un centime, quel est son edge à dix secondes par
séance~?
\end{exercise}

\begin{exercise}[$\star\star$]\label{exo:hf:queue-reactive-models-and-large-tick-assets:4}
Pourquoi une position à l'arrière d'une file peut-elle valoir davantage lorsque la file opposée est longue que lorsqu'elle est courte~?
\end{exercise}

\begin{exercise}[$\star\star$]\label{exo:hf:queue-reactive-models-and-large-tick-assets:5}
Expliquer pourquoi entrer sur une part favorable puis rester a donné de moins bons résultats que toujours entrer.
\end{exercise}

\begin{exercise}[$\star\star$]\label{exo:hf:queue-reactive-models-and-large-tick-assets:6}
D'après la \cref{fig:hf:queue-reactive-models-and-large-tick-assets:sweep}, pourquoi l'edge par séance atteint-il son maximum à un seuil plus bas
que le mark-out par action~?
\end{exercise}

\begin{exercise}[$\star\star\star$]\label{exo:hf:queue-reactive-models-and-large-tick-assets:7}
\emph{Programmation.} Avec le modèle ajusté, calculer la valeur, la probabilité d'exécution et le mouvement adverse d'un lot en tête d'une file acheteuse de 10 lots lorsque
la file vendeuse compte 2 lots (\texttt{order\_value(model(), 0, 10, 2, paths=4000, seed=6)}).
\end{exercise}

\begin{exercise}[$\star\star\star$]\label{exo:hf:queue-reactive-models-and-large-tick-assets:8}
\emph{Trouver la faille.} «~Le \omterm{def:hf:queue-reactive-models-and-large-tick-assets:qr}{modèle réactif à la file} affirme que chaque place dans nos files vaut au moins un dixième de tick, donc nous devrions toujours coter.~»
\end{exercise}

\section{Problème~: en tête de file}

\begin{problem}[{problème du week-end --- en tête de file}]\label{pb:hf:queue-reactive-models-and-large-tick-assets:1}
Un teneur de marché sur une action à un tick doit décider, message après message, dans quelles files se placer.

\medskip\noindent\textbf{Partie I --- Le régime.}
\begin{enumerate}
\item Qu'est-ce qui distingue un actif à gros tick, et que peut contrôler un teneur de marché dans ce cas~?
\item Donner la part à un tick de l'action simulée et la médiane de la meilleure file.
\item Pourquoi le cotateur explorateur du chapitre 3 n'a-t-il obtenu aucune exécution derrière le meilleur cours~?
\end{enumerate}

\medskip\noindent\textbf{Partie II --- Le modèle.}
\begin{enumerate}[resume]
\item Définir le \omterm{def:hf:queue-reactive-models-and-large-tick-assets:qr}{modèle réactif à la file}.
\item Décrire l'estimation et donner la forme des intensités ajustées.
\item Pourquoi les annulations augmentent-elles avec la taille de la file sur ce marché~?
\item Décrire la simulation qui évalue une position dans la file, et son gain.
\end{enumerate}

\medskip\noindent\textbf{Partie III --- La valeur d'une position.}
\begin{enumerate}[resume]
\item Donner les valeurs de la tête et de l'arrière de file pour 10 lots contre 2, et pour 2 lots contre 20.
\item Pourquoi la tête de file ne vaut-elle pas davantage que l'arrière lorsque notre file est courte et l'autre longue~?
\item Pourquoi toutes les valeurs du modèle sont-elles positives, et qu'est-ce que cela omet~?
\end{enumerate}

\medskip\noindent\textbf{Partie IV --- Les règles sur le marché.}
\begin{enumerate}[resume]
\item Définir la \omterm{def:hf:queue-reactive-models-and-large-tick-assets:rules}{règle d'entrée dans la file} et la \omterm{def:hf:queue-reactive-models-and-large-tick-assets:rules}{règle de sortie de la file}.
\item Énoncer le \emph{résultat nommé}~: le mark-out à dix secondes par action de l'entrée systématique, de l'entrée suivie du maintien et de l'entrée suivie de la sortie, ainsi que la
valeur de la tête de file contre l'arrière d'une file de dix lots dans le modèle.
\item Pourquoi la colonne P\&L n'est-elle pas utilisée pour classer les règles~?
\item Pourquoi le maintien des ordres en tête de file a-t-il fait disparaître le gain~?
\item Quel seuil maximise l'edge par séance, et pourquoi pas le mark-out~?
\item Combien de messages chaque règle envoie-t-elle, par rapport à l'entrée systématique, et pourquoi~?
\item Comment ajouter un terme d'information au modèle~?
\item Que changerait-il sur une plateforme pratiquant l'allocation au prorata~?
\item Comment les règles se comporteraient-elles sur un actif à petit tick~?
\item En une phrase~: qu'est-ce qui compte davantage dans un carnet à un tick, bien entrer ou bien sortir~?
\end{enumerate}
\end{problem}

\section{Questions d'entretien}

\begin{interviewq}[$\star$ \iqroles{trader}]\label{iq:hf:queue-reactive-models-and-large-tick-assets:1}
Vous êtes en tête d'une longue file acheteuse sur une action à un tick. La file vendeuse vient de se réduire à un lot. Que faites-vous~?
\end{interviewq}

\begin{interviewq}[$\star\star$ \iqroles{chercheur}]\label{iq:hf:queue-reactive-models-and-large-tick-assets:2}
Comment estimeriez-vous la valeur d'une position en tête de file à partir de vos propres données~?
\end{interviewq}

\begin{interviewq}[$\star\star$ \iqroles{chercheur}]\label{iq:hf:queue-reactive-models-and-large-tick-assets:3}
Pourquoi la probabilité d'exécution et la sélection adverse évoluent-elles ensemble le long d'une file~?
\end{interviewq}

\begin{interviewq}[$\star\star$ \iqroles{développeur}]\label{iq:hf:queue-reactive-models-and-large-tick-assets:4}
Comment un teneur de marché sait-il combien d'actions précèdent son ordre, et qu'est-ce qui fausse l'estimation~?
\end{interviewq}

\begin{interviewq}[$\star\star$ \iqroles{trader}]\label{iq:hf:queue-reactive-models-and-large-tick-assets:5}
Un collègue propose d'annuler puis de remettre vos ordres chaque seconde pour suivre le carnet. Combien cela coûte-t-il sur une action à un tick~?
\end{interviewq}

\begin{interviewq}[$\star\star\star$ \iqroles{chercheur}]\label{iq:hf:queue-reactive-models-and-large-tick-assets:6}
Chaque ordre au repos s'annule au taux $\nu$ et les ordres au marché arrivent au taux $\mu$. Un ordre a $n$ ordres devant lui. Quelle est la probabilité que
le prochain événement parmi ceux qui le précèdent soit un ordre au marché, et quelle est l'espérance du temps jusqu'à ce que les $n$ ordres aient tous disparu~?
\end{interviewq}
